% ============================================================
% Section page (based on ref/starter_page_1.png + ref/starter_page_2.png)
% Índex entry: Fonaments de programació — 3
% ============================================================
\accentfolio
\marginfooter{Programació Didàctica}

\articletitle{Fonaments de programació}

\metabox{Mòdul 1 \;\textbullet\; Durada: 2 setmanes \;\textbullet\; Modalitat: presencial}

\vspace{4mm}
Aquesta secció introdueix els conceptes bàsics de la programació des
d'una mirada didàctica: variables, estructures de control i les
primeres passes en un llenguatge de text o en un entorn de blocs. No
es dona per fet cap coneixement previ i totes les activitats es
plantegen perquè es puguin traslladar directament a l'aula.

\vspace{5mm}
\columndivider{75.5mm}{112mm}{195mm}

\noindent
\begin{minipage}[t]{0.47\textwidth}
\subhead{Què n'aprendràs?}
A identificar les estructures fonamentals de qualsevol llenguatge de
programació i a triar, per a cada grup d'edat, l'eina i el nivell
d'abstracció més adequats.

Es treballen exemples concrets d'aula, des de Scratch fins a
llenguatges de text senzills, sempre partint d'un problema real.
\end{minipage}%
\hfill
\begin{minipage}[t]{0.47\textwidth}
\subhead{Com s'avalua?}
Mitjançant una petita seqüència didàctica dissenyada pel mateix
alumnat, aplicada i revisada amb els companys de mòdul.

\vspace{2mm}
{\small
Lliurament: seqüència didàctica\\
Pes en la nota final: 20\,\%
}
\end{minipage}
